\documentclass[10pt]{article}
% Shared typographic and algebraic source.
\usepackage[letterpaper,margin=0.86in]{geometry}
\usepackage[T1]{fontenc}
\usepackage[utf8]{inputenc}
\usepackage{lmodern}
\usepackage{microtype}
\usepackage{amsmath,amssymb,amsthm,mathtools,mathrsfs}
\usepackage{booktabs,array,tabularx,longtable}
\usepackage{enumitem}
\usepackage[dvipsnames]{xcolor}
\usepackage{fancyhdr,titlesec}
\usepackage[hidelinks]{hyperref}

\definecolor{ink}{HTML}{1F2A33}
\definecolor{teal}{HTML}{226A73}
\definecolor{muted}{HTML}{65747E}
\definecolor{line}{HTML}{C5D1D3}
\definecolor{wash}{HTML}{F3F7F6}

\setlength{\parindent}{0pt}
\setlength{\parskip}{0.46em}
\setlength{\emergencystretch}{2em}
\allowdisplaybreaks[2]
\renewcommand{\arraystretch}{1.16}
\setlist{nosep,leftmargin=1.5em}

\pagestyle{fancy}
\fancyhf{}
\fancyhead[L]{\footnotesize\sffamily\color{muted}COPPALINI ARCHITECTURE}
\fancyhead[R]{\footnotesize\sffamily\color{muted}PHYSICAL REALITY INSIDE CONSTRAINT}
\fancyfoot[C]{\footnotesize\sffamily\color{muted}\thepage}
\renewcommand{\headrulewidth}{0.35pt}
\renewcommand{\headrule}{\hbox to\headwidth{\color{line}\leaders\hrule height \headrulewidth\hfill}}
\setlength{\headheight}{19pt}

\titleformat{\section}{\Large\bfseries\color{ink}}{\thesection.}{0.55em}{}
\titleformat{\subsection}{\large\bfseries\color{teal}}{\thesubsection.}{0.5em}{}
\titlespacing*{\section}{0pt}{1.42em}{0.48em}
\titlespacing*{\subsection}{0pt}{1.02em}{0.32em}

\newtheorem{theorem}{Theorem}
\newtheorem{proposition}[theorem]{Proposition}
\newtheorem{lemma}[theorem]{Lemma}
\theoremstyle{definition}
\newtheorem{definition}[theorem]{Definition}
\theoremstyle{remark}
\newtheorem{remark}[theorem]{Boundary}

\newcommand{\RR}{\mathbb R}
\newcommand{\CC}{\mathbb C}
\newcommand{\QQ}{\mathbb Q}
\newcommand{\ZZ}{\mathbb Z}
\newcommand{\Id}{I}
\newcommand{\tr}{\operatorname{tr}}
\newcommand{\im}{\operatorname{im}}
\newcommand{\kerop}{\operatorname{ker}}
\newcommand{\Span}{\operatorname{span}}
\newcommand{\diag}{\operatorname{diag}}
\newcommand{\snoc}{\operatorname{snoc}}
\newcommand{\Admit}{\operatorname{Admit}}
\newcommand{\Gate}{\operatorname{Gate}}
\newcommand{\Typed}{\operatorname{Typed}}
\newcommand{\Future}{\operatorname{future}}
\newcommand{\Scope}{\operatorname{sc}}
\newcommand{\SilScope}{\mathsf{Silence}_{\rm scope}}
\newcommand{\SilFit}{\mathsf{Silence}_{\rm fit}}

\newcommand{\papertitle}[2]{%
  \begin{center}
  {\color{teal}\rule{\textwidth}{1.35pt}}\\[0.38in]
  {\footnotesize\sffamily\bfseries\color{teal} COPPALINI ARCHITECTURE}\\[0.16in]
  {\fontsize{31}{36}\selectfont\bfseries\color{ink}#1\par}
  \vspace{0.15in}
  {\Large\color{teal}#2\par}
  \vspace{0.22in}
  {\color{line}\rule{\textwidth}{0.55pt}}\\[0.12in]
  {\normalsize Anthony Vito Coppa}\par
  {\small 12 November 2026}\par
  \vspace{0.30in}
  \end{center}
}

\newcommand{\sharedroot}{%
Let
\[
Q=\begin{pmatrix}1&1\\1&0\end{pmatrix},\qquad
\Omega=\begin{pmatrix}0&1\\-1&0\end{pmatrix},\qquad
C=-\Omega,
\]
\[
G=[Q,\Omega]=\begin{pmatrix}-2&1\\1&2\end{pmatrix},\qquad
H=\frac{G}{\sqrt5},\qquad
S=\frac{2Q-I}{\sqrt5}.
\]
The common exact identities are
\[
Q^2=Q+I,\qquad C^2=-I,\qquad G^2=5I,\qquad H^2=S^2=I,
\]
\[
H=CS,\qquad SG=\sqrt5\Omega=-GS,\qquad
HC=-CH,\qquad HCH=-C,
\]
\[
L_s:=e^{sS},\qquad
L_s^{\mathsf T}GL_s=G,\qquad HL_sH=L_{-s}.
\]
}

\newcommand{\sharedoccurrence}[2]{%
A physical field has an interface
\[
\begin{aligned}
\mathfrak I_J:=(&J,\Gamma_J,\operatorname{Src}_J;
A_J,X_J,Y_J,T_J,\Sigma_J,\operatorname{Con}_J,\\
&\rho_J,O_J,K_J,B_J,\theta_J),
\qquad
\rho_J:X_J\to Y_J,\quad
O_J:X_J\to\mathcal O_J,\quad
\theta_J:X_J\to\{0,1\}.
\end{aligned}
\tag{#1}
\]
Its occurrence type is
\[
\begin{aligned}
\mathsf{Occ}_J:=\{&
(\mathfrak I_J;\Theta,x,\rho_J(x),\kappa,
\boldsymbol\varepsilon,\theta_J(x)):\\
&\Theta\in A_J,\ x\in X_J,\ \kappa\in K_J,\
\boldsymbol\varepsilon\in\RR_{\ge0}^{3}\}.
\end{aligned}
\tag{#2}
\]
An indexed occurrence $\mathsf{Occ}_{J,n}$ is an element of
$\mathsf{Occ}_J$.
}

\newcommand{\sharedmatching}{%
Let
\[
\mathcal A_8:=\{\infty\}\sqcup\ZZ_7,
\qquad
(\infty,0,1,2,3,4,5,6)
=(\mathsf{Gr},\mathsf{Li},\mathsf{Me},\mathsf{Ti},
\mathsf{Hz},\mathsf{Si},\mathsf{Sc},\mathsf{Mo}).
\]
For $r\in\ZZ_7$, set
\[
\mathsf M_r=\{\{\infty,r\}\}\cup
\{\{r+j,r-j\}:j=1,2,3\}.
\]
Then $E(K_8)=\bigsqcup_{r\in\ZZ_7}\mathsf M_r$, and the canonical matching is
\[
\mathsf M_0=\{\mathsf{Gravity\! -\! Light},\;
\mathsf{Memory\! -\! Motion},\;
\mathsf{Time\! -\! Scope},\;
\mathsf{Horizon\! -\! Silence}\}.
\]
}

\hypersetup{pdftitle={Suture Field Dynamics: Regular Hamiltonian Reduction},pdfauthor={Anthony Vito Coppa},pdfsubject={Constrained action, regular suture reduction, and Einstein-Maxwell degrees of freedom}}
\newcommand{\hateta}{\widehat\eta}
\newcommand{\hatC}{\widehat C}
\newcommand{\Fsol}{\mathcal F_{\rm sol}}
\newcommand{\rankop}{\operatorname{rank}}
\newcommand{\DoF}{\operatorname{DoF}}
\begin{document}
\papertitle{Suture Field Dynamics}{Regular Hamiltonian reduction}

\begin{abstract}
The suture field supplies an eight-dimensional internal carrier for a
Lorentzian metric and an antisymmetric flux form. This note gives the exact
condition under which its multiplier action has the same propagating content
as Einstein--Maxwell theory. The condition is the rank-sixteen regularity of
the pointwise suture map from the thirty-two coframe coordinates to the ten
metric and six flux coordinates. On that branch the remaining sixteen coframe
coordinates are vertical gauge coordinates; the multiplier sector is
second-class; and the reduced Hamiltonian carries two gravitational and two
electromagnetic polarizations. The reduced theory is Einstein--Maxwell on the
domain where the rest-frame magnetic invariant of the flux is below the
calibration, $b<|\chi|$; the electric invariant is unbounded.
\end{abstract}

\section{Constrained action}

Let $\Pi^A{}_{\mu}$ be a rank-four $E$-valued coframe, with $A=1,\ldots,8$.
Let $B^{\mu\nu}=-B^{\nu\mu}$ and
$\lambda^{\mu\nu}=\lambda^{\nu\mu}$. Fix a nonzero constant calibration
$\chi$. Define
\[
h_{\mu\nu}(\Pi):=\hateta_{AB}\Pi^A{}_{\mu}\Pi^B{}_{\nu},
\qquad
f_{\mu\nu}(\Pi):=\chi\hatC_{AB}\Pi^A{}_{\mu}\Pi^B{}_{\nu}.
\tag{SD.0}
\]
The action is
\[
\begin{aligned}
\mathcal S=\int_M\Bigg[&\frac{c^3}{16\pi G_N}(R(g)-2\Lambda)
-\frac1{4\mu_0}F_{\mu\nu}F^{\mu\nu}\\
&+\frac12B^{\mu\nu}\bigl(F_{\mu\nu}-f_{\mu\nu}(\Pi)\bigr)
+\lambda^{\mu\nu}\bigl(g_{\mu\nu}-h_{\mu\nu}(\Pi)\bigr)\Bigg]
\sqrt{-g}\,d^4x,
\qquad F_{\mu\nu}=2\partial_{[\mu}A_{\nu]}.
\end{aligned}
\tag{SD.1}
\]
Variation of the multipliers gives the two suture conditions
\[
F_{\mu\nu}=f_{\mu\nu}(\Pi),
\qquad g_{\mu\nu}=h_{\mu\nu}(\Pi).
\tag{SD.2}
\]
They are constraint equations. Their multiplier momenta belong to the
canonical system and remain present until Dirac reduction.
In particular, variation with respect to $B$ and $\lambda$ gives SD.2;
it does not eliminate $B$ or $\lambda$. Their values can only be related to
the other equations after the coupled Euler system has been imposed, and
that later relation does not remove their primary momenta from the Dirac
analysis.

\section{The pointwise suture map}

At one spacetime address, define
\[
\mathscr S:\operatorname{Hom}(\RR^4,\RR^8)_{\mathrm{rank}\,4}
\longrightarrow \operatorname{Sym}^2(\RR^4)^*\oplus\Lambda^2(\RR^4)^*,
\qquad
\Pi\longmapsto(h(\Pi),f(\Pi)).
\tag{SD.3}
\]
The domain has dimension $32$ and the target has dimension $10+6=16$. Its
differential is
\[
\begin{aligned}
(D\mathscr S_\Pi\,\delta\Pi)^{(g)}_{\mu\nu}
&=\hateta_{AB}\bigl(\delta\Pi^A{}_{\mu}\Pi^B{}_{\nu}
+\Pi^A{}_{\mu}\delta\Pi^B{}_{\nu}\bigr),\\
(D\mathscr S_\Pi\,\delta\Pi)^{(F)}_{\mu\nu}
&=\chi\hatC_{AB}\bigl(\delta\Pi^A{}_{\mu}\Pi^B{}_{\nu}
+\Pi^A{}_{\mu}\delta\Pi^B{}_{\nu}\bigr).
\end{aligned}
\tag{SD.4}
\]

\begin{definition}[Regular suture branch]
For the fixed $\chi\ne0$, the regular branch is the open locus on which
\[
\rankop D\mathscr S_\Pi=16.
\tag{SD.5}
\]
\end{definition}

An exact certificate that this locus is nonempty is
\[
\Pi_\star=
\begin{pmatrix}
-1&-2& 1& 0\\
 2&-1& 1& 0\\
-1&-2& 1&-2\\
-1& 1&-1& 0\\
-1& 0& 0&-1\\
 0& 1& 1&-1\\
 2& 2& 1&-2\\
-1&-1& 0&-1
\end{pmatrix}.
\tag{SD.6}
\]
For $\chi=1$, exact row reduction gives
\[
\rankop D\mathscr S_{\Pi_\star}=16,
\qquad
h(\Pi_\star)=
\begin{pmatrix}
10&-5& 9& 0\\
-5& 9&-1& 0\\
 9&-1& 6&-5\\
 0& 0&-5&11
\end{pmatrix}.
\tag{SD.7}
\]
Its leading principal minors are
\[
10,\qquad65,\qquad-259,\qquad-4474,
\tag{SD.8}
\]
so the $LDL^{\mathsf T}$ pivot signs are $+,+,-,+$ and
$h(\Pi_\star)$ has Lorentzian signature $(3,1)$. The determinant of a
nonzero $16\times16$ minor of $D\mathscr S$ remains nonzero on a
neighborhood of $\Pi_\star$, which establishes the stated open regular
branch.

\section{Internal vertical directions}

The linear internal stabilizer is
\[
\mathfrak k=
\left\{X\in\mathfrak{gl}(8,\RR):
X^{\mathsf T}\hateta+\hateta X=0,\quad
X^{\mathsf T}\hatC+\hatC X=0\right\}.
\tag{SD.9}
\]
Exact linear reduction gives
\[
\dim\mathfrak k=10,
\qquad
\mathfrak k\cong\mathfrak{so}(1,1)\oplus\mathfrak u(3).
\tag{SD.10}
\]
Every $X\in\mathfrak k$ gives a vertical direction $\delta\Pi=X\Pi$ and
therefore preserves both $h$ and $f$ to first order.  The infinitesimal orbit
map $\alpha_\Pi:\mathfrak k\to\ker D\mathscr S_\Pi$, $X\mapsto X\Pi$, has
point-dependent rank.  It has exact rank $10$ at $\Pi_\star$ and at
$\Pi_{\rm vac}$ (SD.14a), but it is not free on the whole regular branch.
Indeed, with columns
\[
\Pi_{\rm iso}=(\mathbf o_1,\ \mathbf o_2+\mathbf o_3,\ \mathbf o_4,\ \mathbf o_5)
=
\begin{pmatrix}
1&0&0&0\\
0&1&0&0\\
0&1&0&0\\
0&0&1&0\\
0&0&0&1\\
0&0&0&0\\
0&0&0&0\\
0&0&0&0
\end{pmatrix},
\tag{SD.10a}
\]
one has
\[
\rankop D\mathscr S_{\Pi_{\rm iso}}=16,
\qquad
h(\Pi_{\rm iso})=
\begin{pmatrix}-2&1\\1&3\end{pmatrix}\oplus I_2,
\qquad
\rankop\alpha_{\Pi_{\rm iso}}=9.
\tag{SD.10b}
\]
The metric is Lorentzian because the displayed $2\times2$ block has negative
determinant $-7$.  Moreover
$X_{\rm iso}:=0_6\oplus C\in\mathfrak k$ is nonzero and
$X_{\rm iso}\Pi_{\rm iso}=0$, which supplies the isotropy explicitly.

On the regular branch,
\[
\dim\ker D\mathscr S_\Pi=32-16=16.
\tag{SD.11}
\]
Thus the local fiber dimension is constant even though the
$\mathfrak k$-orbit dimension need not be.  Any local complement
$V_\Pi$ with
$\ker D\mathscr S_\Pi=\im\alpha_\Pi\oplus V_\Pi$ has
$\dim V_\Pi=16-\rankop\alpha_\Pi$: it has dimension $6$ at the two rank-ten
certificates and dimension $7$ at $\Pi_{\rm iso}$.  The Hamiltonian reduction
below uses the full sixteen-dimensional vertical fiber, not a constant-rank
$\mathfrak k$-orbit.

\section{Image of the suture map}

For a Lorentzian $h$ and $f\in\Lambda^2(\RR^4)^*$, the endomorphism $h^{-1}f$
is $h$-antisymmetric, so its spectrum is $\{\pm ib,\pm e\}$ with $b,e\ge0$:
\[
b(h,f),\ e(h,f)\ \text{are the rest-frame magnetic and electric invariants of }f.
\tag{SD.11a}
\]
For $f=0$ and for null $f\ne0$, $b=e=0$.

\begin{proposition}[Magnetic cap]
Let $\Pi$ have rank four with $h=h(\Pi)$ Lorentzian, $f=f(\Pi)$, and $\chi\ne0$. Then
\[
(\chi^2-b^2)(\chi^2+e^2)=\chi^4\,\frac{\det\bigl(h+\tfrac{i}{\chi}f\bigr)}{\det h}\ \ge\ 0,
\qquad b\le|\chi|.
\tag{SD.11b}
\]
Equality $b=|\chi|$ holds only at critical points of $\mathscr S$; on the
regular branch $b<|\chi|$.
\end{proposition}

\begin{proof}
Put $F=\chi^{-1}f=\Pi^{\mathsf T}\hatC\Pi$ and $H_c:=\hateta+i\hatC$, a
Hermitian matrix. Then $h+iF=\Pi^\dagger H_c\Pi$. On a Euclidean block
$I_2+iC$ has eigenvalues $0,2$; on the root block $\det(G+iC)=-6$, so it has
signature $(1,1)$. Hence $H_c$ has signature $(4,1)$ with a three-dimensional
kernel, and its restriction to $\im\Pi\otimes\CC$ has at most one negative
eigenvalue. A vector $w\in\RR^4$ with $h(w,w)<0$ gives
$(\Pi w)^\dagger H_c\,\Pi w=h(w,w)<0$, because the $\hatC$ term vanishes by
antisymmetry. So the restriction has exactly one, and $\det(h+iF)\le0$; also
$\det h<0$. For the spectrum $\{\pm ib_F,\pm e_F\}$ of $h^{-1}F$,
\[
(1-b_F^2)(1+e_F^2)=\det(iI-h^{-1}F)=\frac{\det(ih-F)}{\det h}
=\frac{\det(h+iF)}{\det h}\ \ge\ 0,
\]
since $ih-F=i(h+iF)$. With $b=|\chi|b_F$ and $e=|\chi|e_F$ this is SD.11b.
Suppose $b=|\chi|$ at a regular $\Pi$. Then $\mathscr S$ is a submersion
there, so $(h,(1+\varepsilon)f)$ lies in its image for small $\varepsilon$,
with magnetic invariant $(1+\varepsilon)|\chi|>|\chi|$ for $\varepsilon>0$,
contradicting the cap.
\end{proof}

Write $\mathbf o_1,\dots,\mathbf o_8$ for the basis in which
$\hateta=G\oplus I_2\oplus I_2\oplus I_2$ and $\hatC=C\oplus C\oplus C\oplus C$:
$\mathbf o_{1,2}$ span the root plane and $\mathbf o_{3,4}$,
$\mathbf o_{5,6}$, $\mathbf o_{7,8}$ the transverse planes. Let $\pi_\mu$
denote the $\mu$-th column of $\Pi$.

\begin{proposition}[Regular family]
For $a\in\RR$ and $c^2+s^2=1$ let
\[
\pi_0=x\,\mathbf o_1+a\,\mathbf o_7,\quad
\pi_1=\mathbf o_8,\quad
\pi_2=\mathbf o_3,\quad
\pi_3=c\,\mathbf o_4+s\,\mathbf o_5,
\qquad x=\sqrt{\tfrac{1+a^2}2}.
\tag{SD.11c}
\]
Then $h=\diag(-1,1,1,1)$ and
$f=-\chi\bigl(a\,\theta^0\wedge\theta^1+c\,\theta^2\wedge\theta^3\bigr)$, so
$b=|\chi c|$ and $e=|\chi a|$. A $16\times16$ minor of $D\mathscr S$ on a fixed
column set equals $-16\sqrt2\,\chi^{6}s^3(1+a^2)^{7/2}$; the family is regular iff
$s\ne0$, that is, iff $b<|\chi|$.
\end{proposition}

\begin{proposition}[Null witness]
With $x_0=\tfrac{\sqrt6}2$, $x_1=\tfrac{\sqrt6(4-\sqrt2)}{24}$,
$y_1=-\tfrac{\sqrt3}6$, $\beta=\tfrac{\sqrt2}4$, let
\[
\pi_0=x_0\mathbf o_1+\mathbf o_4+\mathbf o_7,\quad
\pi_1=x_1\mathbf o_1+y_1\mathbf o_2+\mathbf o_4+\beta\mathbf o_8,\quad
\pi_2=\mathbf o_3,\quad
\pi_3=\mathbf o_5.
\tag{SD.11d}
\]
Then $h=\eta_4$, $f=\chi(\theta^0+\theta^1)\wedge\theta^2$ (null, nonzero), and
$\rankop D\mathscr S_\Pi=16$ exactly in $\QQ(\sqrt2,\sqrt3)$.
\end{proposition}

\begin{theorem}[Image]
\[
\mathscr S(\text{regular branch})
=\{(h,f):\ h\ \text{Lorentzian},\ b(h,f)<|\chi|\}.
\tag{SD.11e}
\]
\end{theorem}

\begin{proof}
``$\subseteq$'' is SD.11b. For ``$\supseteq$'': for $A\in GL(4)$,
$\mathscr S(\Pi A)=(A^{\mathsf T}hA,A^{\mathsf T}fA)$, and $D\mathscr S_{\Pi A}$
has the rank of $D\mathscr S_\Pi$. By the standard normal forms of a 2-form
relative to a Lorentzian metric, a non-null $f$ is congruent to
$(\eta_4,\ \beta_0\theta^2\wedge\theta^3+\xi_0\theta^0\wedge\theta^1)$
with $|\beta_0|=b$ and $|\xi_0|=e$; a null $f\ne0$ is congruent to
$(\eta_4,\ k(\theta^0+\theta^1)\wedge\theta^2)$, where a boost sets $|k|=|\chi|$
and a reflection its sign; and $f=0$ is congruent to $(\eta_4,0)$. SD.11c with
$a=-\xi_0/\chi$, $c=-\beta_0/\chi$ and $s=\sqrt{1-c^2}>0$ covers the
non-null case, SD.11d the null case, and $\Pi_{\rm vac}$ of SD.14a the case
$f=0$.
\end{proof}

Consequently every solution of SD.1 on the regular branch satisfies
$b(g,F)<|\chi|$ pointwise, every Einstein--Maxwell configuration with
$b(g,F)<|\chi|$ lies in the image, and $e$ is unbounded.

\begin{remark}
The normal forms used in SD.11e are standard and are cited, not proved. The
value of $\chi$ and the units of $b$ are fixed only by calibration: $|\chi|$
must exceed the largest $b$ of the configurations the reduction is to cover.
No empirical comparison is asserted.
\end{remark}

\section{The coframe bracket matrix}

The primary coframe momenta are $P_A{}^\mu\approx0$. Their stabilization
contains
\[
\Psi_A{}^\mu=\sqrt{-g}\left(
2\lambda^{\mu\nu}\hateta_{AB}\Pi^B{}_{\nu}
+\chi B^{\mu\nu}\hatC_{AB}\Pi^B{}_{\nu}\right)\approx0.
\tag{SD.12}
\]
With the canonical convention
$\{\Pi^A{}_{\mu}(x),P_B{}^\nu(y)\}=\delta^A_B\delta_\mu^\nu
\delta^3(x-y)$, the cross-bracket is
\[
\{P_A{}^\mu(x),\Psi_B{}^\nu(y)\}
=-\Delta_{AB}{}^{\mu\nu}(x)\,\delta^3(x-y),
\qquad
\Delta_{AB}{}^{\mu\nu}
=2\sqrt{-g}\left(
\lambda^{\mu\nu}\hateta_{AB}
+\frac\chi2B^{\mu\nu}\hatC_{AB}\right).
\tag{SD.13}
\]
Changing the Poisson-bracket ordering reverses this displayed overall sign;
the rank statement is unchanged.
At the vacuum multiplier point
\[
\lambda^{\mu\nu}=0,\qquad B^{\mu\nu}=0,
\tag{SD.14}
\]
one has $\Delta=0$. Thus nondegeneracy of $\hateta$ and $\hatC$, and rank
four of $\Pi$, do not establish invertibility of $\Delta$. The $32$ coframe
primaries and $32$ displayed stabilizations therefore do not form a globally
invertible second-class block.

This point occurs on a Lorentzian suture vacuum. Let the four columns of
$\Pi_{\rm vac}$ be the first coordinate vector of each $2\times2$ carrier
block. Then
\[
h(\Pi_{\rm vac})=\diag(-2,1,1,1),
\qquad
f(\Pi_{\rm vac})=0.
\tag{SD.14a}
\]
For $\Lambda=0$, the constant metric $g=h(\Pi_{\rm vac})$ and $A=0$ satisfy
the vacuum Einstein--Maxwell equations with the multiplier values in SD.14.

This is not an isolated point. By SD.12,
$\Psi_A{}^\mu=\sqrt{-g}\,m_I\,\partial q^I/\partial\Pi^A{}_\mu$ with $m$ as in
SD.17, and $D\mathscr S_\Pi$ has full row rank on the regular branch. So
$\Psi\approx0$ forces $m=0$, that is
\[
\lambda^{\mu\nu}=B^{\mu\nu}=0,\qquad \Delta=0
\quad\text{on the whole regular constraint surface,}
\tag{SD.14b}
\]
and SD.15a gives $Y^{\mu\nu}=\mu_0^{-1}\mathcal F^{\mu\nu}$ there.

\section{Regular multiplier reduction}

For an unambiguous Legendre transform, introduce an independent antisymmetric
field $\mathcal F_{\mu\nu}$ and an antisymmetric multiplier $Y^{\mu\nu}$:
\[
\begin{aligned}
\mathcal S_{\rm fo}=\int_M\Bigg[&\frac{c^3}{16\pi G_N}(R-2\Lambda)
-\frac1{4\mu_0}\mathcal F_{\mu\nu}\mathcal F^{\mu\nu}\\
&+\frac12Y^{\mu\nu}\bigl(\mathcal F_{\mu\nu}-2\partial_{[\mu}A_{\nu]}\bigr)
+\frac12B^{\mu\nu}\bigl(\mathcal F_{\mu\nu}-f_{\mu\nu}(\Pi)\bigr)\\
&+\lambda^{\mu\nu}\bigl(g_{\mu\nu}-h_{\mu\nu}(\Pi)\bigr)\Bigg]
\sqrt{-g}\,d^4x.
\end{aligned}
\tag{SD.15}
\]
Variation of $Y$ returns $\mathcal F=dA$, so SD.15 is a first-order form of
SD.1. Variation of $\mathcal F$ gives
\[
Y^{\mu\nu}+B^{\mu\nu}=\frac{1}{\mu_0}\mathcal F^{\mu\nu}.
\tag{SD.15a}
\]
Thus $B$ remains an auxiliary multiplier variable through the canonical
reduction; SD.15a is not a license to remove its primary constraint before
that reduction.

By the submersion theorem, local coordinates on the regular coframe carrier
are
\[
\Pi\longleftrightarrow(q^I,\xi^a),
\qquad
q^I=\mathscr S^I(\Pi),
\qquad I,a=1,\ldots,16.
\tag{SD.16}
\]
In these coordinates, the multiplier sector has normal form
\[
\mathcal L_{\rm mult}=\sqrt{-g}\;m_I\bigl(X^I-q^I\bigr),
\qquad
X=(g,\mathcal F),
\qquad
m_{(\mu\nu)}=(2-\delta_{\mu\nu})\lambda^{\mu\nu},\quad
m_{[\mu\nu]}=B^{\mu\nu}\ (\mu<\nu).
\tag{SD.17}
\]
Here $I$ runs over the ten components $\mu\le\nu$ of the symmetric sector and
the six components $\mu<\nu$ of the antisymmetric sector, so that
$\sum_{\mu\nu}\lambda^{\mu\nu}(g-h)_{\mu\nu}
+\frac12\sum_{\mu\nu}B^{\mu\nu}(\mathcal F-f)_{\mu\nu}
=\sum_I m_I(X^I-q^I)$.
Its three Euler relations are
\[
X^I-q^I=0,
\qquad m_I=0,
\qquad \frac{\delta\mathcal L_{\rm mult}}{\delta\xi^a}=0.
\tag{SD.18}
\]
The $\xi^a$ are vertical gauge coordinates. The remaining auxiliary
constraints admit the local Darboux pairing
\[
(P_{qI},q^I-X^I),
\qquad
(P_{mI},m_I),
\tag{SD.19}
\]
which gives $64$ second-class constraints. The $16$ constraints
\[
P_{\xi a}\approx0
\tag{SD.20}
\]
are first class and generate vertical suture motions.

More explicitly, write
\[
\mathcal C_I:=q^I-X^I.
\tag{SD.20a}
\]
\par\noindent
The auxiliary Hamiltonian has the local form
\[
H_{\rm aux}=\int d^3x\,
\bigl(m_I\mathcal C_I
+u_m^IP_{mI}+u_q^IP_{qI}+u_\xi^aP_{\xi a}\bigr).
\tag{SD.20b}
\]
\par\noindent
Its primary stabilizations are
\[
\dot P_{mI}\approx0\Longrightarrow\mathcal C_I\approx0,
\qquad
\dot P_{qI}\approx0\Longrightarrow m_I\approx0,
\qquad
\dot P_{\xi a}\approx0\Longrightarrow0.
\tag{SD.20c}
\]
\par\noindent
With ordered constraint vector
\[
\boldsymbol\zeta=(P_q,\mathcal C;P_m,m),
\tag{SD.20d}
\]
\par\noindent
the second-class matrix is
\[
\{\boldsymbol\zeta(x),\boldsymbol\zeta(y)\}
=
\begin{pmatrix}
0&-I_{16}&0&0\\
I_{16}&0&0&0\\
0&0&0&-I_{16}\\
0&0&I_{16}&0
\end{pmatrix}\delta^3(x-y).
\tag{SD.20e}
\]
\par\noindent
It is invertible. Stabilization of $\mathcal C_I$ and $m_I$ fixes
$u_q^I$ and $u_m^I$; it creates no tertiary auxiliary constraint.

\par\noindent
\emph{Lapse and shift.} Because $N$ and $N^i$ are components of $g$, the
primary constraints $\pi_a=(\pi_N,\pi_i)$ have nonzero brackets with the
$g_{0\mu}$-components of $\mathcal C_I$:
$\{\pi_a,\mathcal C_J\}=\partial X^J/\partial n^a$ with $n^a=(N,N^i)$. The
combinations
\[
\tilde\pi_a:=\pi_a+\frac{\partial X^I}{\partial n^a}P_{qI}
\tag{SD.20f}
\]
commute with every $\mathcal C_J$ and every $P_{qJ}$. The bracket matrix of
$(\pi_a,P_q,\mathcal C;P_m,m)$ therefore has rank $64$ and a four-dimensional
kernel spanned by the $\tilde\pi_a$, which replace $\pi_N,\pi_i$ in SD.24,
SD.27 and SD.29. Write $K_a{}^I:=\partial X^I/\partial n^a$.  Since
$\mathcal H_\perp$ and $\mathcal H_i$ contain neither $n^a$ nor the auxiliary
coordinates $q^I$, one has
\[
\{\tilde\pi_a,\mathcal H_A\}
=\{K_a{}^I,\mathcal H_A\}P_{qI}\approx0,
\qquad A\in\{\perp,i\}.
\tag{SD.20g}
\]
Thus the modified lapse--shift primaries are first class modulo the
second-class surface.  Likewise, by definition of the submersion chart
$q=\mathscr S(\Pi)$, neither $X$, $\mathcal C=q-X$, nor the physical
Hamiltonian depends on $\xi$; hence
\[
\{P_{\xi a},H_{\rm phys}\}=0,
\qquad
\{P_{\xi a},\mathcal C_I\}=0,
\qquad
\{P_{\xi a},\mathcal H_A\}=0
\tag{SD.20h}
\]
strongly.  The sixteen $P_{\xi a}$ therefore generate the full local vertical
gauge fiber independently of the pointwise rank of the linear stabilizer
orbit.

\par\noindent
\emph{Charts.} The pairing SD.19 is written for SD.15, where $\mathcal F$ is an
independent field: phase-space dimension $148$, with $26$ first-class and $88$
second-class constraints, and $\DoF=\frac12(148-2\cdot26-88)=4$. For the
direct action SD.1 counted in SD.28, $F_{0i}$ is a velocity, and the same
reduction holds with $F_{0i}$ expressed through the momentum of $A_i$ and
$B^{0i}$: primary $5+32+6+10=53$, secondary $5+6+10+16=37$, total
$90=26+64$.

\begin{theorem}[Regular suture reduction]
On the regular branch, Dirac reduction of SD.15 is locally equivalent, on the
domain $b(g,F)<|\chi|$ of SD.11e, to
\[
\mathcal S_{\rm red}=\int_M\left[
\frac{c^3}{16\pi G_N}(R-2\Lambda)
-\frac1{4\mu_0}F_{\mu\nu}F^{\mu\nu}
\right]\sqrt{-g}\,d^4x,
\qquad F=dA.
\tag{SD.21}
\]
\par\noindent
\end{theorem}

\begin{proof}
Let $H_{\rm phys}$ denote the Einstein--Maxwell first-order Hamiltonian part
of SD.15, including its spatial-derivative terms, and let the auxiliary part
be SD.20b.  With the ordering SD.20d, the inverse of the second-class kernel
SD.20e is
\[
C^{-1}(x,y)=
\begin{pmatrix}
0&I_{16}&0&0\\
-I_{16}&0&0&0\\
0&0&0&I_{16}\\
0&0&-I_{16}&0
\end{pmatrix}\delta^3(x-y).
\tag{SD.21a}
\]
Accordingly, preservation of $\mathcal C_I$ gives
\[
\dot{\mathcal C}_I
=\{\mathcal C_I,H_{\rm phys}\}+u_q^I\approx0,
\qquad
u_q^I=-\{\mathcal C_I,H_{\rm phys}\},
\tag{SD.21b}
\]
while preservation of $m_I$ fixes $u_m^I$.  Thus any spatial derivatives in
$H_{\rm phys}$ enter the value of the second-class multiplier; they do not
supply an additional constraint.

For any two physical functionals $A,B$ independent of
$(q,P_q,m,P_m,\xi,P_\xi)$, the Dirac bracket is
\[
\{A,B\}_D=\{A,B\}
-\iint d^3z\,d^3w\,
\{A,\zeta_\alpha(z)\}C^{-1}_{\alpha\beta}(z,w)
\{\zeta_\beta(w),B\}.
\tag{SD.21c}
\]
Such functionals may have nonzero brackets with
$\mathcal C=q-X$, but they commute with $P_q,P_m,m$.  In SD.21a the
$\mathcal C\mathcal C$ block is zero and each $\mathcal C$ entry is paired
only with $P_q$.  Therefore every correction term in SD.21c vanishes:
\[
\boxed{\{A,B\}_D=\{A,B\}}.
\tag{SD.21d}
\]
Hence eliminating
$P_q=0$, $\mathcal C=0$, $P_m=0$, and $m=0$ leaves the physical Poisson
algebra unchanged.  Equation SD.20h quotients the remaining $\xi$ directions;
no physical term depends on them.  Variation of $Y$ gives
$\mathcal F=dA$, and imposing $q=X$ then removes the auxiliary copy of the
suture coordinates.  Substitution yields exactly SD.21 on the image domain
SD.11e.  Finally SD.20g shows that the lapse--shift primaries retain their
first-class role after the second-class reduction.
\end{proof}

\section{ADM Hamiltonian and stabilization}

Let
\[
K_{ij}=\frac1{2N}\bigl(\dot\gamma_{ij}-D_iN_j-D_jN_i\bigr),
\qquad
\pi^{ij}=\frac{c^3\sqrt\gamma}{16\pi G_N}
\bigl(K^{ij}-\gamma^{ij}K\bigr),
\tag{SD.22}
\]
\par\noindent
and, after the multiplier reduction, let
\[
E^i=-\frac{\sqrt{-g}}{\mu_0}F^{0i}.
\tag{SD.23}
\]
\par\noindent
The reduced primary constraints are
\[
\pi_N\approx0,\qquad \pi_i\approx0,\qquad E^0\approx0.
\tag{SD.24}
\]
\par\noindent
The total Hamiltonian is
\[
H_T=\int d^3x\left(
N\mathcal H_\perp+N^i\mathcal H_i+A_0\mathcal G
+u_N\pi_N+u^i\pi_i+u_0E^0\right),
\tag{SD.25}
\]
\par\noindent
where
\[
\begin{aligned}
\mathcal H_\perp={}&
\frac{16\pi G_N}{c^3\sqrt\gamma}
\left(\pi^{ij}\pi_{ij}-\frac12\pi^2\right)
-\frac{c^3\sqrt\gamma}{16\pi G_N}\bigl({}^{(3)}R-2\Lambda\bigr)\\
&+\frac{\mu_0}{2\sqrt\gamma}\gamma_{ij}E^iE^j
+\frac{\sqrt\gamma}{4\mu_0}F_{ij}F^{ij},\\
\mathcal H_i={}&-2\gamma_{ik}D_j\pi^{jk}+F_{ij}E^j,
\qquad
\mathcal G=-\partial_iE^i.
\end{aligned}
\tag{SD.26}
\]
\par\noindent
Preservation of SD.24 gives
\[
\dot\pi_N\approx0\Longrightarrow\mathcal H_\perp\approx0,
\qquad
\dot\pi_i\approx0\Longrightarrow\mathcal H_i\approx0,
\qquad
\dot E^0\approx0\Longrightarrow\mathcal G\approx0.
\tag{SD.27}
\]
\par\noindent
The scalar, momentum, and Gauss constraints close under the standard
hypersurface-deformation and $U(1)$ gauge algebra. Their stabilizations fix
only the corresponding Dirac multipliers and produce no tertiary physical
constraint.

\section{Degree count}

On the regular Darboux chart of the direct action SD.1, the phase-space
dimensions per spatial address are:
\[
\begin{array}{c|c}
\text{sector}&\text{phase-space dimension}\\ \hline
(\gamma_{ij},N,N^i) &20\\
A_\mu &8\\
\Pi^A{}_{\mu}&64\\
B^{\mu\nu}&12\\
\lambda^{\mu\nu}&20\\ \hline
\text{total}&124
\end{array}
\tag{SD.28}
\]
There are ten Einstein--Maxwell first-class constraints and sixteen vertical
suture first-class constraints. The regular multiplier sector supplies sixty
four second-class constraints. Therefore
\[
\DoF=
\frac12\bigl(124-2(10+16)-64\bigr)=4.
\tag{SD.29}
\]
Equivalently, after exact reduction the Einstein--Maxwell phase space has
dimension $28$ and ten first-class constraints:
\[
\DoF_{\rm red}=\frac12\bigl(28-2\cdot10\bigr)=4.
\tag{SD.30}
\]
The four propagating degrees of freedom are the two gravitational and two
electromagnetic transverse polarizations.

\section{Discrete carrier}

The $\ZZ_7$ Fourier action of the suture carrier preserves $\hateta$ and
$\hatC$. In SD.1 it is an algebraic fiber symmetry and contributes no new
canonical pair. A local connection or a measured holonomy coupling enters
only through an additional specified dynamical term and calibrated
occurrence record.

\end{document}
